% Begin


%%%%%%%%%%%%%%%%%%%%%%%%%%%%%%%%%%%%%%%%%%%%%%%%%%%%%%%%%%%%%%%
\chapter*{\S \space 49. --- HOMOMORPHISMES DE $\mathcal{M}_{0, 3}$, LES GROUPES $\mathcal{M}_{\rho,\sigma}$ GÉNÉRAUX}\thispagestyle{empty}
\addcontentsline{toc}{section}{{\bf 49}. Les groupes $\mathcal{M}_{\rho,\sigma}$ généraux}
\label{sec:49}
\section*{}

%%%%%%%%%%%%%%%%%%%%%%%%%%%%%%%%%%%%%
\addcontentsline{toc}{subsection}{I. Préliminaires heuristiques}
\vskip .3cm
{\bf I)} Considérons\footnote{Où on se met dans le bain\dots Toute cette section I semble entièrement inutile pour la suite, c'est un tâtonnement préliminaire qui débouche sur la section II, p.\ 553.} un sous-groupe fermé $\mathcal{G}$ de $\mathcal{M}_{0,3}$, contenant $\mathfrak{S}^{+\wedge}_{0, 3}$, donc image inverse d'un sous-groupe fermé de ${\boldsymbol{\Gamma}_{0,3}}^\sim$. On suppose que celui-ci est ``suffisament gros'' --- disons qu'il contienne un sous-groupe d'indice fini de $\boldsymbol{\Gamma}_{\mathbf{Q}} \subset {\boldsymbol{\Gamma}_{0,3}}^\sim$. Pour le moment, on suppose pour simplifier que $\mathcal{G}\subset \Ker \varepsilon_3$. En résumé
$$
\mathfrak{S}_{0, 3}^{+ \wedge} \subset \mathcal{G} \subset \mathcal{M}_{0, 3}^\sim, \quad u \in \mathcal{G} \Rightarrow \mu(u) \equiv 1 (3).
\leqno (1)
$$
Cette dernière hypothèse signifie donc que pour tout $u \in \mathcal{G}$, on a $\alpha_0(u)\in\mathfrak{S}_{0, 3}^{\wedge +}$, donc on a 
$$
\mathcal{G} \overset{\underline{\alpha_0}}{\underset{\underline{\beta_0}}{\rightrightarrows}} \mathfrak{S}_{0,3}^{+},
\leqno (2)
$$
caractérisés par
$$
(2') \quad
\begin{cases}
u(\rho) = \operatorname{int}(\alpha_0)\rho\;, \quad \alpha_0 = \underline{\alpha_0}(u)\; , \\
u(\sigma) = \operatorname{int}(\beta_0)\sigma\;, \quad \beta_0 = \underline{\beta_0}(u)\; .
\end{cases}
$$
On considère
\[
\begin{tikzpicture}[x=1pt,y=1pt]

% top line
\node at (0,150) [anchor=west] 
{$\mathcal{G}(0)= \mathcal{G}\cap\mathcal{M}_{0,3}^\sim(0)$};

% bottom line
\node at (0,0) [anchor=west] 
{$\mathcal{G}[0]= \mathcal{G}\cap\mathcal{M}_{0,3}^\sim[0]
= \operatorname{Norm}_{\mathcal{G}}(L_0^{\mathfrak{S}})\; ,$};

% vertical segment
\draw (40,70) -- (40,92);

% top half oval (the "cap")
\draw (40,70) ellipse [x radius=10, y radius=25];

\end{tikzpicture}
\]
de sorte que 
$$
\mathcal{G}= \mathcal{G}(0)\cdot \mathfrak{S}_{0, 3}^{+ \wedge} \quad \text{(produit semidirect)}
\leqno (4)
$$
$$
\mathcal{G}[0]= \mathcal{G}(0)\cdot L_0^\mathfrak{S}\quad\quad \text{idem}\quad
\leqno (5)
$$
$$
\mathcal{G}[0]\cap\mathfrak{S}_{0, 3}^{+ \wedge}= L_0^\mathfrak{S}\; .
\leqno (6)
$$
On a posé
$$
L_0^\mathfrak{S}= \text{sous-groupe fermé de $\mathfrak{S}_{0, 3}^{+ \wedge}$ engendré par $\varepsilon_0\;$,}
\leqno (7)
$$
et on définit de même $L_1^\mathfrak{S}$, $L_\infty^\mathfrak{S}$.

Considérons un groupe profini $\mathcal{M}$, et un homomorphisme (continu) \emph{surjectif}
$$
\varphi\; :\; \mathcal{G}\;\to\;\mathcal{M}\; .
\leqno (8)
$$
On pose
$$
\begin{cases}
\varphi(\mathfrak{S}_{0, 3}^{+ \wedge}) = \mathfrak{S} \quad \text{sous-groupe distingué de $\mathcal{M}$} \\
\varphi(\mathcal{G}(0)) = \mathcal{M}(0) \\
\varphi(\mathcal{G}[0]) = \mathcal{M}[0] \\
\varphi(L_0^\mathfrak{S}) = L_0\; .
\end{cases}
\leqno (9)
$$
Donc ce sont des sous-groupes fermés de $\mathcal{M}$, satisfaisant
$$
\mathcal{M}(0)\; \subset \; \mathcal{M}[0]\; \subset\; \operatorname{Norm}_\mathcal{M}(L_0)
\leqno (10)
$$
$$
L_0\;\subset\; \mathfrak{S}\cap\mathcal{M}[0]\; .
\leqno (11)
$$
Par abus de langage, on désigne encore par $\rho$, $\sigma$, $\varepsilon_0$, $\varepsilon_1$ \dots les images de ces éléments dans $\mathfrak{S}$, qui y satisfont donc les relations connues dans $\mathfrak{S}_{0, 3}^{+ \wedge}$, en particulier
$$
\sigma^2 = \rho^3 = 1\; ,\quad  \sigma\rho= \varepsilon_0\; ,\quad \rho\sigma= \varepsilon_1\; ,\quad\; \varepsilon_1= \sigma(\varepsilon_0) = \rho(\varepsilon_0)\; .
\leqno (12)
$$
Les restrictions des applications $\alpha_0$, $\beta_0$ de (2, 2') à $\mathcal{M}_{\rho, \sigma}[0]$ seront plutôt notées $\alpha,\beta$, où $\alpha$, $\beta$ sont en fait des applications
$$
S\mathcal{G}\; \defeq \; \underbrace{\text{image inverse de $\mathcal{G}$ dans $S\mathcal{M}_{0,3}^\sim$}}_{= \{ (u,f)\in\mathcal{G}\times\mathfrak{S}_{0, 3}^{+ \wedge}\; |\; fu\in\mathcal{G}[0]\}}
    \; \overset{\alpha}{\underset{\beta}{\rightrightarrows}} \mathfrak{S}_{0, 3}^{+ \wedge}
\leqno (13)
$$
dont on regarde les restrictions au sous-groupe ($\isom\mathcal{M}_{\rho, \sigma}[0]$) de $S\mathcal{G}$ défini par la relation $f = 1$.

Composant $\underline{\alpha},\underline{\beta}:\mathcal{G}[0]\rightrightarrows\mathfrak{S}_{0, 3}^{+ \wedge}$ avec $\varphi|\mathfrak{S}_{0, 3}^{+ \wedge}$, on trouve des applications, encore notées $\underline{\alpha}$, $\underline{\beta}$ (ou $\underline{\alpha}^\mathcal{M}$, $\underline{\beta}^\mathcal{M}$ en cas de confusion possible)
$$
\mathcal{G}[0] \;\overset{\underline{\alpha}}{\underset{\underline{\beta}}{\rightrightarrows}}\;\mathfrak{S}\; ,
\leqno (14) 
$$
satisfaisant donc
$$
\begin{cases}
u(\rho) = \operatorname{int}(\alpha)\rho \\
u(\sigma) = \operatorname{int}(\beta)\sigma
\end{cases}
\text{\shortstack[l]{pour $u\in\mathcal{M}[0]$ de la forme $\varphi(u_0)$,\\et $\alpha = \alpha(u_0)$, $\beta = \beta(u_0)\;$.\vspace*{-4mm}}}
\leqno (15)
$$
A priori, $\alpha$ et $\beta$ dépendent du choix de $u_0$ dont $u$ provient, pas seulement de $u$ --- mais
$$
(16) \quad
\begin{cases}
\alpha\rho(\alpha)^{-1} = [u,\rho] \\
\beta\sigma(\beta)^{-1} = [u,\sigma]
\end{cases}
$$
n'en dépendent pas. Si on pose
$$
T_\rho=\underline{\operatorname{Centr}}_\mathcal{M}(\rho)\; ,\quad T_\sigma= \underline{\operatorname{Centr}}_\mathcal{M}(\sigma)\; ,
\leqno (17)
$$
on sait donc que, modulo $T_\rho$ resp.\ $T_\sigma$, $\alpha$, $\beta$ ne dépendent pas de $u$.

Hypothèse. --- $\alpha = \alpha(u_0)$, $\beta = \beta(u_0)$ ne dépendent que de $u = \varphi(u_0)$ (pour $u_0\in\mathcal{G}[0]$).

On va exprimer cette hypothèse de fa\c{c}on un peu différente, en introduisant
$$
J = \Ker(\mathcal{G}[0]\to\mathcal{M}[0])\; ,
\leqno (18)
$$
on veut donc que pour
\[
u_0'= \delta_0 u_0\; ,
\]
avec $u_0\in\mathcal{G}[0]$, $\delta_0\in J$, on ait
\[
\alpha(u_0')=\alpha(u_0)\; ,\quad \beta(u_0')=\beta(u_0)\; .
\]
Soient $\alpha_0 = \underline{\alpha}(u_0)$, $\beta_0 = \underline{\beta}(u_0)$, $\alpha'_0 = \underline{\alpha}(u'_0)$, $\beta'_0 = \underline{\beta}(u'_0)$, 
$\overline{\alpha}_0 = \underline{\alpha}(\delta_0)$, $\overline{\beta}_0 = \underline{\beta}(\delta_0)$ \emph{dans} $\mathcal{G}$, de sorte qu'on 
aura, par les formules de cocycles dans $\mathcal{G}$,
\[
\alpha_0'= \delta_0(\alpha_0)\overline{\alpha}_0\; ,\quad \beta_0'= \delta_0(\beta_0)\overline{\beta}_0\; ,
\]
d'où, en appliquant $\varphi$ et notant que
$$
\varphi(\delta_0(\alpha_0)) = \underbrace{\varphi(\delta_0)}_{=\; 1}(\varphi(\alpha_0)) = \varphi(\alpha_0)\;\defeq \; \alpha\; , \quad \text{et de même}
$$
$$
\varphi(\delta_0(\beta_0)) = \varphi(\beta_0)\;\defeq \; \beta
$$
$$
\alpha'= \alpha\overline{\alpha}\; ,\quad \beta'= \beta\overline{\beta}\quad\text{dans $\mathfrak{S}$}\; ,
\leqno (19)
$$
où $\alpha$, $\beta$, $\alpha'$, $\beta'$, $\overline{\alpha}$, $\overline{\beta}$ sont les images par $\varphi$ des quantités $\alpha_0$ etc.\ dans $\mathfrak{S}^{+\wedge}_{0,3}$. Donc l'hypothèse signifie aussi
$$
\text{$\forall\;\delta_0 \in J \defeq \Ker(\mathcal{G}[0]\to\mathcal{M}[0])$, on a $\alpha^\pi(\delta_0),\beta^\pi(\delta_0) \in \Ker \varphi$.}
\leqno (20)
$$
Quand cette hypothèse est satisfaite, on a donc des applications
$$
\mathcal{M}[0]\;\overset{\underline{\alpha}}{\underset{\underline{\beta}}{\rightrightarrows}}\;\mathfrak{S}
\leqno (21)
$$
factorisant (14), et satisfaisant donc
$$
\begin{cases}
u(\rho) = \operatorname{int}(\alpha)\rho \\
u(\sigma) = \operatorname{int}(\beta)\sigma
\end{cases}
\quad\text{pour $u\in\mathcal{M}[0]$, $\alpha = \underline{\alpha}(u)$, $\beta = \underline{\beta}(u)$}\; ,
\leqno (22)
$$
qui paraphrasent (21). 

D'ailleurs, posant
$$
\xi = \alpha \rho (\alpha)^{-1}\; ,\quad \eta= \beta\sigma(\beta)^{-1}\; ,
\leqno (23)
$$
les relations (22) s'écrivent aussi
$$
u(\rho) = \xi\rho\; ,\quad u(\sigma) = \eta\sigma
\leqno (24)
$$
et on aura
$$
\xi = [u,\rho]\; , \quad \eta = [u,\sigma]\; ,
\leqno (25)
$$
et de plus
$$
\xi = \eta \ell_1^\nu \; , \quad \text{où $\nu = \frac{\mu - 1}{2}\in \hat{\mathbf{Z}}\;\;$ ($\mu$ le multiplicateur de $u$)}\; ,
\leqno (26)
$$
qui provient de la relation analogue dans $\mathfrak{S}_{0,3}^+$. (On suppose pour fixer les idées que l'homomorphisme
$$
\mathcal{G}\;\xrightarrow{\chi^\mathcal{G}}\; \hat{\mathbf{Z}}^* \quad\text{(`multiplicateur')}
$$
se factorise par
$$
\mathcal{G}\;\to\;\mathcal{M}\;\xrightarrow{\chi^\mathcal{M}}\; \hat{\mathbf{Z}}^* \; ,
\leqno (27)
$$
de sorte qu'on a une notion de multiplicateur ($\in\hat{\mathbf{Z}}^\ast$) pour les $u\in\mathcal{M}$.)

Considérons l'homomorphisme naturel
$$
\begin{aligned}
\mathcal{M}[0] &\to \Aut(\mathfrak{S}) \\
u &\longmapsto \operatorname{int}(u)|\mathfrak{S}\; ,
\end{aligned}
\leqno{(28)}
$$
on voit sur (22) ou sur (24) qu'il est connu quand on connaît les fonctions $\underline{\alpha},\; \underline{\beta}\; :\; \mathcal{M}[0] \to \mathfrak{S}$, et même quand on connaît seulement $\xi,\; \eta \; :\; \mathcal{M}[0] \to \mathfrak{S}$\footnote{car $\rho$, $\sigma$ engendrent topologiquement $\mathfrak{S}$.}. Si d'ailleurs on pose
$$
\mathcal{M}'[0] = \operatorname{Im}(\mathcal{M}[0]\to\Aut(\mathfrak{S}))\; ,
\leqno (29)
$$
ces dernières fonctions $\xi$, $\eta$ sont déjà définies sur $\mathcal{M}'[0]$ --- alors qu'il n'est pas clair qu'il en soit de même de $\alpha$, $\beta$, quand $\mathcal{M}[0]$ n'opère pas fidèlement sur $\mathfrak{S}$. On a ainsi une application $(\underline{\xi},\underline{\eta})$ d'un certain sous-groupe $\mathcal{M}'[0] \subset \operatorname{Norm}_{\Aut(\mathfrak{S})}(L_0)$ dans l'ensemble des solutions de l'équation en $\xi$, $\eta$ (26) qu'on écrira simplement
$$
\eta^{-1}\xi\;\in\; L_1,
\leqno (31)
$$
relation qui implique (26) (avec la forme précise de l'exposant $\nu$) si on suppose p.ex.\ $\mathfrak{S}$ assez gros pour que $\mathfrak{S}_{0,3}^{+\,\wedge}\to \SL(2,\hat{\mathbf{Z}})$ se factorise par $\mathfrak{S}$ --- ce qu'on va supposer, pour nous simplifier la vie.

De façon plus précise, soit
$$
\Sigma = \Sigma_{\mathfrak{S},{\rho,\sigma}}
= \left\{ (\xi,\eta)\in\mathfrak{S}\times\mathfrak{S}\;\middle|\;
\begin{cases}
\text{a) } \xi\rho \text{ conjugué à } \rho \text{ dans } \mathfrak{S}, \text{ i.e. } \xi = [\alpha,\rho] = \alpha\rho(\alpha)^{-1},\\
\text{b) } \xi\sigma \text{ conjugué à } \sigma \text{ dans } \mathfrak{S}, \text{ i.e. } \eta = [\beta,\sigma] = \beta\sigma(\beta)^{-1},\\
\text{c) } \eta^{-1}\xi\in L_1, \text{ i.e. } \exists\,\nu'\in\hat{\mathbf{Z}} \text{ tel que } \xi = \eta\varepsilon_1^{\nu'}.
\end{cases}
\right\}.
$$
et soit $u\in\Aut(\mathfrak{S})$ un automorphisme de $\mathfrak{S}$, et posons
$$
\left\{
\text{
\begin{tabular}{ll}
$\xi  = [u,\rho]   = u \rho u^{-1} \rho^{-1} = u(\rho)\rho^{-1}\; ,$ & i.e.\ $u(\rho) = \xi\rho$ \\
$\eta = [u,\sigma] = u(\sigma)\sigma^{-1}\; ,$                            & i.e.\ $u(\sigma) = \eta\sigma$ \\
\end{tabular}
}
\right. \; , \;\;\;\text{cf.\ (24), (25),} \;\;\;\text{\tt [\footnotemark]}\; , \hspace*{-5mm}
\leqno (34)
$$
\footnotetext{\tt [(33) n'existe pas.]}%
alors dire que $\xi\rho$ ($= u(\rho)$) est conjugué à $\rho$ signifie que $u$ fixe la classe de conjugaison de $\rho$, dire que $\eta\sigma$ ($= u(\sigma)$)
est conjugué de $\sigma$ signifie que $u$ fixe la classe de conjugaison de $\sigma$. Enfin la condition $\eta^{-1}\xi\in L_1$ s'écrit aussi, comme
\[
\begin{aligned}
\eta^{-1}\xi 
& = & (\sigma u(\sigma^{-1}))(u(\rho)\rho^{-1}) \\
& = & \sigma u(\sigma^{-1}\rho)\rho^{-1} \\
& = & \sigma \Big(u(\underbrace{\sigma^{-1}\rho}_{=\;\varepsilon_0}) \underbrace{(\rho^{-1}\sigma)}_{=\; \varepsilon_0^{-1}} \Big) \\
& = & \sigma(u(\varepsilon_0)\varepsilon_0^{-1})\; ,
\end{aligned}
\quad (\footnotemark)
\]
\footnotetext{NB $\sigma^{-1} = \sigma$.}%
sous la forme
\[
\sigma(u(\varepsilon_0)\varepsilon_0^{-1})\;\in\; L_1\; , \quad\text{i.e.\ $\; u(\varepsilon_0)\varepsilon_0^{-1}\;\in\; L_0$}
\]
(puisque $L_0 = \sigma^{-1}(L_1)$), ou encore (puisque $\varepsilon_0\in L_0$) $u(\varepsilon_0)\in L_0$, soit (puisque $\varepsilon_0$ engendre $L_0$)
\[
u(L_0)= L_0\; ,
\]
i.e.\ que $u$ \emph{normalise} $L_0$ (\footnote{Dans cet argument, on n'a utilisé que la relation $\sigma_0^{-1}\rho = \varepsilon_0$ {\tt [plutôt 
$\sigma^{-1}\rho = \varepsilon_0$]} et le fait que $\varepsilon_0$ engendre $L_0$, et pas que $\sigma^2 = 1$ ou $\rho^3 = 1$ p.ex.}).

Soit donc
$$
\begin{aligned}
\tilde{B}_0' 
& \defeq & \left\{ u\in\Aut(\mathfrak{S})\; \left|\; 
\text{
\begin{tabular}{l}
$u$ fixe {\tt [la]} classe de conjugaison de $\rho$ et celle de $\sigma$ \\ 
$u$ normalise $L_0$ \\
\end{tabular}
} \right.\right\} \\
& \subset & \Aut(\mathfrak{S})\; ,\\
\end{aligned}\hspace*{-5mm}
\leqno (35)
$$
de sorte qu'on a une application (injective parce que $\rho$, $\sigma$ engendrent $\mathfrak{S}$)
$$
\begin{aligned}  
\tilde{B}_0' & \hookrightarrow     & \Sigma \\
u       & \longmapsto & \Big( [u,\rho] = u(\rho)\rho^{-1}\; , \; [u,\sigma] = u(\sigma)\sigma^{-1}\Big)\; .
\end{aligned}
\leqno (36)
$$
$$
\text{
\begin{tabular}{p{14cm}}
\emph{Hypothèse fondamentale sur $\mathfrak{S}$, muni de $\rho$, $\sigma$, d'où $\varepsilon_0 = \sigma^{-1}\rho$, 
$\varepsilon_1 = \rho\sigma^{-1} = \sigma(\varepsilon_0) = \rho(\varepsilon_0)$~:} \\
\begin{tabular}{ll}
a) & $\mathfrak{S}$ est engendré topologiquement par $\rho$, $\sigma$. \\
b) & L'application (36) (injective par a)) est bijective. \\
\end{tabular} \\
\end{tabular}
}
\leqno (37)
$$
La condition b) signifie donc que pour $(\alpha,\beta)\in\mathfrak{S}\times\mathfrak{S}$, satisfaisant une relation
\[
\underbrace{\alpha\rho(\alpha)^{-1}}_{=\;\xi}= \underbrace{\beta\sigma(\beta)^{-1}}_{=\;\eta}\varepsilon_1^{\nu'}\; ,\quad\text{$\nu'$ convenable,}
\]
il existe $u\in\Aut(\mathfrak{S})$ tel qu'on ait
\[
\begin{aligned}
{[u,\rho]} & = & \alpha\rho(\alpha)^{-1}\; , & \text{i.e.} & u(\rho) & = & \operatorname{int}(\alpha)(\rho)\; , & \text{i.e.} & u(\rho) & = & \xi\rho \\
{[u,\sigma]} & = & \beta\sigma(\beta)^{-1}\; , & \text{i.e.} & u(\sigma) & = & \operatorname{int}(\beta)(\sigma)\; , & \text{i.e.} & u(\sigma) & = & \eta\sigma \\
\end{aligned}
\]
(on aura donc nécessairement $u\in\tilde{B}_0'$, et $u$ sera unique si on a a)).

Ainsi, sur l'ensemble des $R_{\rho, \sigma}$ des $(\xi,\eta)$ satisfaisant a) b) c), on aura une structure de groupe, déduite par transport de structure de celle de $\tilde{B}_0'$ etc.


%End
